\documentclass[12pt]{article}
\usepackage[utf8]{inputenc}
\usepackage[T1]{fontenc}
\usepackage{amsmath,amssymb}
\usepackage{geometry}
\geometry{margin=2.5cm}

\begin{document}

\noindent\underline{\textbf{The Fundamental Theorem of Arithmetic}}

\bigskip
\noindent $(\forall)\,a\in\mathbb{N}$, $a>1$, $\Rightarrow a=p_1\cdots p_n$

\noindent [is represented uniquely as a product of prime numbers]

\bigskip
\noindent\underline{Lemma}: $(\forall)\,a\in\mathbb{N}$, $a>1$ \ $\exists\,p$ prime, $p\mid a$.

\bigskip
\noindent $P=\{d \mid d\mid a \text{ and } d\neq 1\}$

\bigskip
\noindent Let $p\in P$ be the smallest number in $P$.

\noindent\underline{$p$ is prime}.

\bigskip
\noindent Suppose, for contradiction, $p\neq$prime $\ \exists\,d_1$ such that $d_1\mid p$, $d_1\neq 1$ and $d_1\neq p$.

\[
\left.\begin{aligned} d_1&\mid p\\ p&\mid a\end{aligned}\right\} \Rightarrow d_1\mid a
\]
\[
\left.\begin{aligned} &\Rightarrow d_1\in P\\ &d_1<p\end{aligned}\right\} \Rightarrow p\text{ is prime.} \quad \times
\]

\bigskip
\noindent\underline{Proof of the theorem}:

\noindent Suppose it is not true.

\noindent Denote $P=\big\{$the set of natural numbers (greater than 1)\footnote{the abbreviation in the original (``not.\ core'') is unclear -- reconstructed here as a ``$>1$'' condition, matching the statement of the theorem.} that cannot be represented as

\noindent a finite product of prime numbers$\big\} \neq \varnothing$.

\bigskip
\noindent Let $a$ be the smallest number in $P$.

\noindent $\Rightarrow a$ is not prime

\bigskip
\noindent Let $a=bc$.

\noindent $b,c\notin P$

\[
\Rightarrow b=p_1\cdots p_j
\]
\[
c=p_1'\cdots p_i'
\]
\[
\Rightarrow a = p_1\cdots p_j\cdot p_1'\cdots p_i'.
\]

\bigskip
\noindent\underline{Uniqueness}: \ $a = p_1\cdots p_n = p_1'\cdots p_m'$.

\[
\overset{(5)}{\Longrightarrow} p_1\mid p_1'\cdots p_m' \Rightarrow \exists\, i \text{ such that } p_1\mid p_i' \Rightarrow p_1=p_i'
\]

\end{document}
